\documentclass[preview]{standalone}

\usepackage{amsmath}
\usepackage{amssymb}
\usepackage{stellar}
\usepackage{definitions}

\begin{document}

\id{finite-abelian-groups-classification}
\genpage

\section{The classification theorem}

\begin{snippettheorem}{finite-abelian-group-elementary-divisor-classification}{Classification by elementary divisors}
    Every finite \abeliangroup \(G\) is a finite direct product of
    \cyclicgroup[cyclic groups] of prime-power order:
    \[
        G\groupisomorphic
        C_{p_1^{\alpha_1}}\times\cdots\times
        C_{p_t^{\alpha_t}},
        \qquad
        \alpha_i\geq1.
    \]
    The multiset of prime powers
    \(\{p_1^{\alpha_1},\ldots,p_t^{\alpha_t}\}\) is uniquely determined
    by \(G\), up to reordering.
\end{snippettheorem}

\begin{snippetproof}{finite-abelian-group-elementary-divisor-classification-proof}{finite-abelian-group-elementary-divisor-classification}{Classification by elementary divisors}
    An \abeliangroup is precisely a \(\integers\)-\module. Since
    \(\integers\) is a \principalidealdomain, the primary decomposition
    theorem applies. Finiteness forces the free rank to be zero, because
    every nonzero free \(\integers\)-module is infinite. Moreover,
    \[
        \integers/(p^\alpha)\groupisomorphic C_{p^\alpha}.
    \]
    Existence and uniqueness now follow directly from the corresponding
    statements for \elementarydivisor[elementary divisors] of finitely
    generated modules over a PID.
\end{snippetproof}

\begin{snippetcorollary}{finite-abelian-group-invariant-factor-classification}{Classification by invariant factors}
    Every finite \abeliangroup \(G\) admits a decomposition
    \[
        G\groupisomorphic
        C_{d_1}\times\cdots\times C_{d_s},
        \qquad
        1<d_1\divides d_2\divides\cdots\divides d_s.
    \]
    The integers \(d_1,\ldots,d_s\) are uniquely determined by \(G\).
\end{snippetcorollary}

\begin{snippetproof}{finite-abelian-group-invariant-factor-classification-proof}{finite-abelian-group-invariant-factor-classification}{Classification by invariant factors}
    Apply invariant-factor decomposition to \(G\) as a
    \(\integers\)-\module. Again, finiteness eliminates the free part,
    and \(\integers/(d_i)\groupisomorphic C_{d_i}\). Choosing each
    \(d_i>1\) removes the unit ambiguity in the
    \invariantfactor[invariant factors].
\end{snippetproof}

\begin{snippetnote}{finite-abelian-groups-fixed-order}{Groups of a fixed finite order}
    If
    \[
        |G|=\prod_{j=1}^{k}p_j^{n_j},
    \]
    then the possible \(p_j\)-primary components are in bijection with
    the partitions of \(n_j\). Consequently, the isomorphism classes of
    finite abelian groups of this order are obtained by choosing
    independently one partition of every exponent \(n_j\).
\end{snippetnote}

\section{When a product of two cyclic groups is cyclic}

\begin{snippetcorollary}{product-two-finite-cyclic-groups-cyclic}{Product of two finite cyclic groups}
    For positive integers \(m,n\),
    \[
        C_m\times C_n\groupisomorphic C_{mn}
        \ifandonlyif
        \gcd(m,n)=1.
    \]
\end{snippetcorollary}

\begin{snippetproof}{product-two-finite-cyclic-groups-cyclic-proof}{product-two-finite-cyclic-groups-cyclic}{Product of two finite cyclic groups}
    If \(\gcd(m,n)=1\), the ideals \((m)\) and \((n)\) of
    \(\integers\) are comaximal. The Chinese remainder theorem gives
    \[
        C_m\times C_n
        \groupisomorphic
        \integers/m\integers\oplus\integers/n\integers
        \groupisomorphic
        \integers/mn\integers
        \groupisomorphic C_{mn}.
    \]

    Conversely, suppose a prime \(p\) divides both \(m\) and \(n\).
    Write
    \[
        m=p^hm',
        \qquad
        n=p^kn',
        \qquad
        h,k\geq1,
        \qquad
        p\nmid m'n'.
    \]
    The \(p\)-primary component of \(C_m\times C_n\) has two
    \elementarydivisor[elementary divisors] \(p^h\) and \(p^k\), whereas
    the \(p\)-primary component of \(C_{mn}\) has the single elementary
    divisor \(p^{h+k}\). Uniqueness of elementary divisors shows that the
    groups are not isomorphic.
\end{snippetproof}

\begin{snippetnote}{two-cyclic-groups-gcd-lcm-invariant-factors}{The gcd--lcm invariant factors}
    More generally,
    \[
        C_m\times C_n
        \groupisomorphic
        C_{\gcd(m,n)}\times C_{\lcm(m,n)}.
    \]
    For every prime \(p\), the two invariant-factor exponents are
    \[
        \min\{v_p(m),v_p(n)\}
        \quad\text{and}\quad
        \max\{v_p(m),v_p(n)\}.
    \]
    Thus the first factor is trivial exactly when \(\gcd(m,n)=1\), which
    is precisely the cyclic case.
\end{snippetnote}

\end{document}
